\section{Compatibility on the actual built family and the full local gain}
\label{sec:witness-realization-v049}

\noindent\textbf{Scope and conclusion.}
The five-field compatibility record constructed in GCF12 and transported
in GCF16 has to concern the comparison map in the supplied arithmetic
object.  An independently chosen compatible comparison does not suffice.
This chapter obtains the required identification directly from the
original built-point class.  Its construction record already contains
one joint equality of the group ties, comparison, twisted forms and
kernel.  The final point changes the division lattice and its height;
we prove that the same joint equality survives that change.  This gives
the original \(\mathrm{GroupCompatibilityFamily}\) on every member of
the original built class, before imposing the selected-design conditions.
Together with the selected tops already treated in GCF18 and the full
local estimates of LGS13, it supplies the five original
\(\mathrm{LinkG}\) conclusions with their original class quantifiers.

The argument does not construct a nonempty design class.  Nor does it
construct the stronger \(\mathrm{RecT}\), its row pin, host record or
Weil-equivariance field.  Those objects are required elsewhere in the
proof.  They are unnecessary for this direct compatibility supplier.
An arbitrary family outside the built class still needs an actual
compatibility input.  A separate final statement also proves the
universal original \texttt{GlobalTheory.Top.link\_g} implication, at its
literal \(\mathrm{TSel}\) and duty premises.

\paragraph{Objects and notation.}
The original \(D\) contains the number field \(L=D.\mathrm{LL}\), its
specified complex embedding, the host, words and places, the column
spaces, their integral models and comparison domains.  A full object
\(O:\mathrm{Obj}(D)\) includes its size data, actual group instance,
local lattices and \(O.\mathrm{ex}:\mathrm{Extra}\).
The smaller object \(\mathrm{ofObj}(O):\mathrm{Obj2}(D)\) retains
\[
 \begin{gathered}
 G_t,\ \mathrm{th},\ c,\ \mathrm{tw},\ G_y,\
 \mathrm{errX},\ \mathrm{nch},\ \mathrm{xd},\ \mathrm{lev},\
 X,\ p_0,\ \mathrm{ker},\ \mathrm{amb},\ \mathrm{amb}_0,\
 \mathrm{gauge}.
 \end{gathered}
\]
In particular it retains the actual \(\mathrm{Extra}.\mathrm{cmp}\)
and \(\mathrm{Extra}.\mathrm{ker}\); forgetting size fields does not
forget either of them.  A prime on an object or map below means the
original \texttt{setDiv} transport.  It never denotes an independent
choice of a local model or a change of coefficient field.

% BEGIN WRS1
\paragraph{WRS1. Division transport preserves the whole local kernel record.}
\label{claim:WRS1}
Let \(N_u\) be arbitrary \(\mathcal O_L\)-submodules of the original
division vector spaces and let \(h_u\in\mathbb R\) be arbitrary.
The original record
\(D'=\mathrm{setDiv}(D,N,h)\) replaces
\[
 \mathrm{Nu}(u).L\ \text{by}\ N_u,\qquad
 \mathrm{div}.h\ \text{by}\ h,\qquad
 \mathrm{HH}\ \text{by}\
 \max\{1,\mathrm{sFrame}(\mathrm{Nu}[L:=N],h)+\mathrm{RR}.\mu_{\max}\}.
 \tag{WRS.1}
\]
All other fields are retained.  For every supplied
\[
 k_{u,w}:\mathrm{LocalToGlobal}_2
       (D,G_t,c,\mathrm{tw},u,w;G_y.\mathrm{Ch}(u,w)),
\]
there is the original transported record on
\((D',G_t',c',\mathrm{tw}',G_y')\), with precisely the same
coefficient embeddings, local lattice maps, class maps, Taylor series
and comparisons.  This holds for every allowed \(u,w\) and the
original good-index guard.  There is no condition that the new
\(N_u\) be full, finitely generated or nonzero.

Here is a field-by-field verification.  Neither the residue prime, the
local Frobenius period, the chain length nor the local filling and level
data changes under (WRS.1).  The two identities identifying these
lengths are therefore unchanged.  Keep the original maps
\(\lambda:\overline{\mathbb Q}\to\mathbb C_p\),
\(i_K:K_f\to\mathbb C_p\), and all
\(i_{w,j,v}:E\to Q_{j,v}\).
Their coefficient-embedding identity remains
\[
 i_K(c.\mathrm{embW}(z))=\lambda(\mathrm{emb}_L z),
 \tag{WRS.2}
\]
with the same algebraicity guard on the right.
Keep the original global-lattice map, its local-span identity, its
class map, the equivalence relating equality of classes to the local
level action, and its surjectivity onto every class at every
\(j\le m\).  Those statements use the twisted column lattices and
local level tuples; they do not use the division submodules replaced
in (WRS.1).

Retain the original Taylor series at every adjacent pair of local
classes.  Finally keep every semilinear comparison \(k.\mathrm{cmp}_j\),
its injectivity for \(j\le m\), and the three exact laws
\[
 \begin{aligned}
 k.\mathrm{cmp}_0(\mathrm{tw}.e_0 x)&=c.\mathrm{cmp}(x),\\
 \mathcal O_{K_f}\operatorname{span}
   \bigl(k.\mathrm{cmp}_0(\mathrm{tw}.\mathrm{lat}_0)\bigr)
     &=\mathrm{Lv.Int}_0,\\
 \mathrm{emb}_{j+1}y=\mathcal K_j(\mathrm{emb}_j x)
   &\ \Longrightarrow\
 k.\mathrm{cmp}_{j+1}y=\mathrm{Ch}.K_{z,j}(k.\mathrm{cmp}_j x)
       \qquad(j<m).
 \end{aligned}                                                    \tag{WRS.3}
\]
All expressions in (WRS.2)--(WRS.3) are the same expressions on the
same vectors before and after transport.  This checks all seventeen
original fields of \(\mathrm{LocalToGlobal}_2\), including the last
chain-compatibility field, not just the underlying vector-space map.
Taking these records for all \(u,w\) gives the original
\texttt{kerTr}.  Transporting a full \(\mathrm{Extra}\) or an
\(\mathrm{Obj2}\) uses this very record at its kernel field.
In particular no auxiliary zero-chain substitution is used here.
\hfill\(\square\)
% END WRS1

% BEGIN WRS2
\paragraph{WRS2. Taylor presentation and every kernel tie survive division.}
\label{claim:WRS2}
Fix original true-column data \(w,Z\), a column record
\(t:\mathrm{ColT}(w,Z)\), a divisor record \(X\), a local lattice
record \(G_y\), and
\[
 k:\mathrm{KerFam}(w.DZ,t.G_t,t.c,w.\mathrm{tw}(Z),G_y),\qquad
 \mathcal K:\mathrm{KerT}(t,X,G_y,k).
\]
For every \(N,h\) as in WRS1 there is a full original
\[
 \mathcal K':\mathrm{KerT}
       (t',X',G_y',k'),                              \tag{WRS.4}
\]
where the column transport is the one of GCF16 and the local kernel
transport is the one of WRS1.  If \(\mathcal K\) satisfies the
original \(\mathrm{EntC}\) predicate, then \(\mathcal K'\) does too,
with the same place embeddings and the same basis matrices.

The nine fields of \(\mathrm{KerT}\) can all be retained.  The
lifting objects \(\mathcal O_{u,\mu}\) have the same Hermitian
spaces, frames, theta functions, lattices and transformation rules.
The standardness assertion uses \(X.Q_p\), which is retained by
\(X'=\mathrm{divTr}(X)\).  All the unfolding assertions at
\(i<m\) therefore concern the same functions.  Keep the global
kernel presentation \(E\).  Its compatibility \(\mathrm{EntOk}\)
with the lifting objects keeps the same witnesses at each tuple,
place and level; the transported objects have those same fields.
The equality \(k_{u,w}.\lambda=E.\lambda_w\), the complete served
Taylor-series clause, the symmetric kernel embedding equation, and
the positive-place equation retain identical maps on identical
arguments.  This supplies, respectively, the fields
\(\mathrm{O},\mathrm{std},\mathrm{unf},\mathrm{E},\mathrm{ok},
\mathrm{lam},\mathrm{served},\mathrm{iw},\mathrm{iw\_pos}\).

For clarity, preserving the presentation in \(\mathrm{EntC}\) is
stronger than finding some isomorphic presentation.  Write its given
witness as
\[
 E=\mathrm{kerEntC}(P,\mathcal O,B;H),
 \qquad H:\mathrm{KerEntHyp}(P,\mathcal O,B,X,\mathrm{tw}).
 \tag{WRS.5}
\]
The polynomial weight spaces, evaluation maps and their kernels do
not involve the three changed fields in (WRS.1).  Hence their quotients
\(V_{u,\mu,i}=\mathrm{polV}_{u,i}/\ker(\mathrm{evP}_{u,\mu,i})\)
are literally the same quotients.  Retain their sections, the
coefficient embeddings \(P.i_L\), all matrices \(B_{u,\mu,w,i,c,\tau}\),
and the inverse maps used on the original polynomial image.  This
retains the source and target functionals \(\mathrm{prS}\),
\(\mathrm{prT}\), the Taylor series \(\mathrm{Th}\), and the
normalizing scalar \(\mathrm{unC}\).  Its convention replacing a
zero raw scalar by \(1\) is also unchanged.

The four clauses of \(H\) now remain the same: the target
coefficient norm is the norm of a coefficient-field element; the
source lattice is characterized by all source minor functionals
having norm at most one; the target lattice has the corresponding
target characterization; and the kernel formula is
\[
 \mathrm{prT}(v',I)=\mathrm{unC}\,
 \mathrm{prS}\!\left(v,
       \mathrm{applyAt0}(\mathrm{Th},\mathrm{Gxz}(I))\right)
 \tag{WRS.6}
\]
whenever the given global kernel maps \(\mathrm{evC}(v)\) to
\(\mathrm{evC}(v')\).  The original level-free guard also stays
unchanged, since it reads \(X\) and the twisted column data, not
the replaced division lattice.  Thus the very same \(H\) supplies
\(H'\), and every field in the definition of \(\mathrm{kerEntC}\)
is the corresponding transported field in (WRS.5).
Applying this transport to that one equality proves
\(E'=\mathrm{kerEntC}(P',\mathcal O',B;H')\).
This proves the full original
\texttt{kernelTaylorSeries\_preserved\_divisionLatticeChange},
with its \(\mathrm{EntC}\) conclusion when supplied originally.
\hfill\(\square\)
% END WRS2

% BEGIN WRS3
\paragraph{WRS3. The dependent comparison--kernel equality transports once.}
\label{claim:WRS3}
Suppose \(O_2:\mathrm{Obj2}(w.DZ)\), a column record \(t\), and
\(k_0\) satisfy the original equality
\[
 (O_2.G_t,O_2.c,O_2.\mathrm{tw},O_2.\mathrm{ker})
    =(t.G_t,t.c,w.\mathrm{tw}(Z),k_0).                \tag{WRS.7}
\]
This is an equality in the dependent sum whose fourth entry has
the type
\[
              \mathrm{KerFam}(D,G_t,c,\mathrm{tw},O_2.G_y).
\]
For every \(N,h\), the corresponding transported equality is
\[
 (O_2'.G_t,O_2'.c,O_2'.\mathrm{tw},O_2'.\mathrm{ker})
    =(t'.G_t,t'.c,w.\mathrm{tw}(Z'),k_0').            \tag{WRS.8}
\]
Here \(Z'=\mathrm{setDivZ}(Z,N,h)\), so
\(w.DZ'=\mathrm{setDiv}(w.DZ,N,h)\).

Define a map on the entire dependent sum in (WRS.7): apply
\(\mathrm{gatTr}\), \(\mathrm{cmpTr}\) and
\(\mathrm{twTr}\) to its first three entries and apply WRS1 to
its fourth entry, using the transported \(O_2.G_y\).  WRS1 shows
that the fourth entry has exactly the type required by the first
three transported entries.  Thus this is a map on the whole sum,
not four independent identifications of records of different types.
The image of the left side is the left side of (WRS.8), by the
definition of \(\mathrm{Obj2.tr}\).  GCF16 identifies the
transported column's first entries with the specified
\(\mathrm{gatTr},\mathrm{cmpTr}\), and its twisted forms with
\(w.\mathrm{tw}(Z')\).  The right image is therefore exactly the
right side of (WRS.8).  Equality is preserved by this one map,
which proves the assertion.  If \(\mathrm{KerT}\) is supplied,
WRS2 transports it on this same \(k_0'\).
\hfill\(\square\)
% END WRS3

% BEGIN WRS4
\paragraph{WRS4. A built core supplies all moved column ties on the given object.}
\label{claim:WRS4}
Let \(b:\mathrm{Built}(\mathcal S,K,w,Z,c_0)\) be the original
built-core witness, for arbitrary original slots \(\mathcal S\).
Let \(N,h\) be arbitrary division data and let \(O\) be any full
object over \(D'=\mathrm{setDiv}(w.DZ,N,h)\) such that
\[
             \mathrm{ofObj}(O)=c_0.O.\mathrm{tr}(N,h).
 \tag{WRS.9}
\]
There is the specified transported
\(t'=\mathrm{columnLocalDataTransportDivision}(b.\mathrm{col})\)
such that
\[
 \begin{gathered}
 \mathrm{ColsExtraTiedS}(t'.\mathrm{full},O.\mathrm{ex}),\\
 \mathcal G:\mathrm{GroupCompatibilityAll}
             (D',O.G_t,O.\mathrm{ex}.c),\\
 \mathcal G=t'.g,\qquad
 \mathcal G.\mathrm{herm}=t'.\mathrm{full}.\mathrm{to0},\qquad
 \mathcal G.\mathrm{forms}(u)=t'.C.\mathrm{forms}(u).\mathrm{to0}.
 \end{gathered}                                                   \tag{WRS.10}
\]
As usual the equality of records in (WRS.10) is read after the
joint dependent identification of their indexing data.  The given
kernel of \(O.\mathrm{ex}\) is part of this conclusion.

The \(\mathrm{Built}\) witness contains its actual
\(b.\mathrm{col}:\mathrm{ColumnDataSplitPrimes}\), its kernel
family \(k_0\), its \(\mathrm{KerT}\), and precisely (WRS.7).
It also contains
\[
 c_0.O.\mathrm{lev}=b.\mathrm{col.full.base.lev}_0,
 \qquad c_0.O.\mathrm{errX}=\mathrm{toErrXS}(b.\mathrm{col.full}).
 \tag{WRS.11}
\]
GCF16 provides \(t'\) with the original coefficient field, the
full five compatibility fields, all column readings, \(\mathrm{IsRecG}\)
and \(\mathrm{OmegaCompatible}\).  WRS2 provides the transported
kernel tie, and WRS3 gives the joint identity (WRS.8).
Substitute that identity once and then use (WRS.9).

The three fields of \(\mathrm{ColsExtraTiedS}\), expanded as
\(\mathrm{SymmetricExtraCompatibility}\), are now available on
the actual \(O.\mathrm{ex}\).  Its join field consists of the two
equalities (WRS.11) and the column record's exact prime-compatibility
law.  Its kernel-embedding field is the transported
\(b.\mathrm{ker.iw}\), on the given kernel, and its positive-place
field is \(b.\mathrm{ker.iw\_pos}\) on that same kernel.  Explicitly,
the former preserves at every allowed \(u,w,v\) the identity
\[
 e_{Q,v}\bigl((O.\mathrm{ex.ker})_{u,w}.i_w(0,v,z)\bigr)
   =\operatorname{pr}_1\bigl(\mathrm{spl}_v(\mathrm{incl}_E z)\bigr)
       \qquad(z\in E).                              \tag{WRS.12}
\]
Use \(t'.g\) itself for \(\mathcal G\).  Its Hermitian and form
equalities are already part of the transported column record.
These prove every assertion in (WRS.10).  Neither \(\mathrm{EntC}\),
a new row choice, a full \(\mathrm{Wit}\), nor a selector is a
premise of this statement.  The original built-core witness is its
actual input; it already supplies the column and kernel data used.
\hfill\(\square\)
% END WRS4

% BEGIN WRS5
\paragraph{WRS5. The original built point has that exact object.}
\label{claim:WRS5}
Let \(\beta:\mathrm{BPt}(K)\) satisfy the original
\(\beta.\mathrm{Is}(\mathrm{IsBuilt}(\mathcal S))\).
Then (WRS.10), including its full symmetric extra compatibility,
holds on the actual \((\beta.D,\beta.O)\).
The compatibility record is on \(\beta.O.\mathrm{ex}.c\), and
its kernel equations hold on \(\beta.O.\mathrm{ex.ker}\).

By the definition of \(\mathrm{IsBuilt}\), take the original
\(w,Z,c_0,b\) with the single joint equality of the columns,
later stage and core to \(\beta.X,\beta.c\).  Substitute this
dependent equality; it identifies the base with \(D=w.DZ\)
and the core with \(c_0\).  The remaining data of \(\beta\)
contain a layout \(Y\) and its specified equality
\[
                 \mathrm{layObj}(Y)=c_0.O,
 \tag{WRS.13}
\]
together with its given \(p:\mathrm{LedPart}(K.C,Y)\).

The source's final point is not an unspecified isomorphic model.
Its division lattice and height are exactly
\[
 \begin{aligned}
 N'_u&=\mathcal O_L\operatorname{span}
       \{a y:a\in\mathrm{Iu}(D,p.L,u),\ y\in p.L.\mathrm{NuLat}(u)\},\\
 h'_u&=\mathrm{hExpl}_6(D,p.L,Y.X,K.C.\mathrm{Csz})(u),\\
 p.\mathrm{pt}.D&=\mathrm{setDiv}(D,N',h').
 \end{aligned}                                                    \tag{WRS.14}
\]
Before this operation, the full object \(Y.\mathrm{obj}(p)\)
uses the given size and group data of \(p\), and the comparison,
twisted forms, local lattices and kernel of \(Y\).
Thus its \(\mathrm{Obj2}\) is precisely \(\mathrm{layObj}(Y)\).
The definition of \(p.\mathrm{pt}\) transports that full object
by (WRS.14).  Reading each field therefore gives
\[
 \begin{aligned}
 \mathrm{Obj.ofPt}(p.\mathrm{pt})
       &=\bigl(Y.\mathrm{obj}(p)\bigr).\mathrm{tr}(N',h'),\\
 \mathrm{ofObj}(\beta.O)
       &=\mathrm{layObj}(Y).\mathrm{tr}(N',h')
        =c_0.O.\mathrm{tr}(N',h').
 \end{aligned}                                                    \tag{WRS.15}
\]
For the local lattice field, the point uses \(Y.G_y\); the
original \(p.G_y\)-equality transports the group tie to that same
field.  Hence (WRS.15) identifies the entire kernel domain as well
as its comparison, rather than merely forgetting that distinction.

Apply WRS4 with exactly \(N',h',O=\beta.O\).
No size bound on (WRS.14) is needed for this application, and none
of \(p\)'s original sizes, metrics, prime sets or constants is
replaced.  The output is the full original compatibility on the
actual point.  In particular its coefficient field is the one
already in \(\beta.D\).
\hfill\(\square\)
% END WRS5

% BEGIN WRS6
\paragraph{WRS6. Every point of the original built class receives the same-data supplier.}
\label{claim:WRS6}
Let \(Q\) be any original core predicate refining
\(\mathrm{IsBuilt}(\mathcal S)\).  Keep arbitrary original bundle
type \(\mathrm{Bd}\), linkage predicate \(\mathrm{Lk}\), and
further property \(E\).  Define the original class
\[
                \mathcal A=\mathrm{Fc.cls}(Q,\mathrm{Lk},E).
 \tag{WRS.16}
\]
For every \(\mathrm{Fm}\in\mathcal A\), every
\(x>\mathrm{Fm}.x_1\), every witness \(h_x\) of that inequality,
and every \(i\in\mathrm{Fm}.I_x\), the actual object
\[
 O_{x,h_x,i}=
 (\mathrm{Fm}.L,\mathrm{Fm}.G_t,\mathrm{Fm}.G_x,
  \mathrm{Fm}.G_y,\mathrm{Fm}.G_K,\mathrm{Fm}.X,
  \mathrm{Fm}.c_X,\mathrm{Fm}.\mathrm{ex})_{x,h_x,i}
 \tag{WRS.17}
\]
has the full conclusion (WRS.10).  No additional \(\mathrm{TSel}\)
or eventual witness-realization hypothesis is needed.

Indeed class membership supplies one original parameter record
\(\kappa\), its \(\mathrm{Link}(\kappa,\mathrm{Fm})\), and
the seven-duty class whose object property is
\(\mathrm{QOf}(Q,\mathrm{Lk},E,\kappa)\).
The link supplies the required positivity of \(\kappa\).
Apply this object property at the specified \(x,h_x,i\), using
the family's actual \(\mathrm{Csz}\).  The definition supplies
a bundle \(b\) and an actual built witness
\(K,\beta,b'\), together with
\[
  \bigl(\mathrm{Fm}.D(x,i),O_{x,h_x,i},b\bigr)
       =\bigl(\beta.D,\beta.O,b'\bigr).              \tag{WRS.18}
\]
This is the original equality in the dependent sum of the base,
full object and bundle, not a collection of numerical agreements.
Also \(\beta.\mathrm{Is}(Q)\) holds.  Its refinement gives the
built-core witness required by WRS5.  Apply that claim and substitute
(WRS.18) once.  It identifies the coefficient rings, the comparison
map, the complete kernel family, and their indexing records with
the ones in (WRS.17).  This proves the assertion for the originally
specified point and hence for all of them.

The proof does not change \(E\) to insert compatibility and does
not change \(\mathrm{Lk}\).  Compatibility follows from the
\emph{existing} built witness in their class definition.  For the
original selected predicate \(Q=\mathrm{IsGoodE}(\mathcal S)\),
the required refinement is immediate from its definition: forget
\(\mathrm{GoodO}\) and \(\mathrm{EntC}\), retaining its actual
\(\mathrm{Built}\) witness.  The same applies to any further
conjunction used in the original common class.
\hfill\(\square\)
% END WRS6

% BEGIN WRS7
\paragraph{WRS7. Full family compatibility and its original Closed order.}
\label{claim:WRS7}
For the class \(\mathcal A\) of WRS6 and every subclass
\(\mathcal C\subseteq\mathcal A\),
\[
 \begin{gathered}
 \mathcal C(\mathrm{Fm})\Longrightarrow
                 \mathrm{GroupCompatibilityFamily}(\mathrm{Fm}),\\
 \mathrm{Closed}\bigl(\mathcal C,
           (\mathrm{Fm},C_H)\mapsto
                    \mathrm{GroupCompatibilityFamily}(\mathrm{Fm})\bigr).
 \end{gathered}                                                   \tag{WRS.19}
\]
In particular these hold on \(D_3(3,\mathcal A)\), on its
target-compatible subclass, and on the original \(W_3(\mathcal A)\).

The first predicate in (WRS.19) is exactly
\[
 \begin{gathered}
 \exists x_g\ \forall x\ge x_g\ \forall i\ \forall h_x:
          \mathrm{Fm}.x_1<x,\\
 \operatorname{Nonempty}\!\left(
 \mathrm{GroupCompatibilityAll}
       (\mathrm{Fm}.D(x,i),\mathrm{Fm}.G_t(x,i),
                       \mathrm{Fm}.\mathrm{ex}(x,h_x,i).c)\right).
 \end{gathered}
 \tag{WRS.20}
\]
Take \(x_g=\mathrm{Fm}.x_1+1\).  WRS6 applies to every later
\(x,i,h_x\), on the exact \(\mathrm{Extra}\) named in (WRS.20).
This uniform threshold precedes all places, words, local tuples
and coefficient completions.  It also precedes all choices of
\(h_x\): the source class property is universal in that argument,
and each application of (WRS.18) uses its specified value.

To check the second assertion, expand the source's \(\mathrm{Closed}\):
\[
 \begin{gathered}
 \exists\epsilon_1>0\ \forall\epsilon\in(0,\epsilon_1]\
 \exists K_0:\mathrm{TransportedAlt}\\
 \forall\mathrm{Fm}:\mathrm{FamP}(\epsilon)\ \forall C_H\in\mathbb R,
 \qquad \mathcal C(\mathrm{Fm})\ \land\
       \mathrm{Fm.Large}(K_0,C_H)\Longrightarrow P(\mathrm{Fm},C_H).
 \end{gathered}
 \tag{WRS.21}
\]
Choose \(\epsilon_1=1\) and choose the bound with all four
components zero.  This choice is made before \(\mathrm{Fm}\)
and \(C_H\).  The first assertion of (WRS.19) proves the
conclusion without using the large-coefficient guard; consequently
it certainly proves (WRS.21) with that guard.  The family threshold
\(x_g\) is chosen afterwards, in its permitted position.
This argument is valid on every fixed subclass, so it also applies
after adding the original smoothness, twisting, selected-top,
\(\mathrm{TuSize}\), \(\mathrm{DGamSize}\) and target predicates.
None of those operations changes the members' objects.
\hfill\(\square\)
% END WRS7

% BEGIN WRS8
\paragraph{WRS8. The second Closed input uses the original selected tops.}
\label{claim:WRS8}
For the very same \(\mathcal A\), the original selected class
\[
 D_3(3,\mathcal A)=\mathrm{Des}(\mathrm{WithChN}(\mathcal A),3)
 \tag{WRS.22}
\]
satisfies
\[
       \mathrm{Closed}(D_3(3,\mathcal A),\mathrm{Adm2F}).
 \tag{WRS.23}
\]
Its \(\mathrm{ChosenE}\) field supplies \(\mathrm{TopsF}\)
on the actual chosen subset.  In particular this is the same
family, selected subset and degree assignment as in WRS7.

Expand membership in (WRS.22).  It includes
\(\exists C_H'\,\mathrm{ChosenE}(\mathrm{Fm},3,C_H')\)
in addition to the original base-class, smoothness and twisting
conditions and \(\mathrm{TuSize},\mathrm{DGamSize}\).
GCF18 projects its \(\mathrm{ChosenTT}\) and one positive
\(C_b\) with \(\mathrm{ChosenBad}\).  The completed estimate
GLG16 then gives \(\mathrm{Adm2F}\) with the original threshold
\[
 x_{\mathrm{top}}=
 \max\{x_a,x_b,\mathrm{Fm}.x_1+1,
                         (16C_b+2)^{3/\epsilon}\}.
 \tag{WRS.24}
\]
This is a theorem about the supplied selected tops, not a new
choice of a random subset or a later restriction of the family.
It holds for every member of the class before imposing any
coefficient bounds.  The proof of (WRS.21), again using cutoff
\(1\) and zero coefficient bounds, therefore proves (WRS.23).
The internal witness \(C_H'\) in \(\mathrm{ChosenE}\) need
not equal the external \(C_H\) universally quantified in
\(\mathrm{Closed}\); neither this projection nor its estimate
asserts such an equality.  All thresholds retain their permitted
dependence on the fixed family and its chosen-top witnesses.
\hfill\(\square\)
% END WRS8

% BEGIN WRS9
\paragraph{WRS9. The full original LinkG on the selected built class.}
\label{claim:WRS9}
For every original built class \(\mathcal A\) in WRS6,
\[
             \mathrm{Closed}(D_3(3,\mathcal A),\mathrm{LinkG}).
 \tag{WRS.25}
\]
This has no additional \(\mathrm{TSel}\), \(\mathrm{WPt}\),
kernel-existence or compatibility premise.  The actual compatibility
input is supplied by WRS7 and the actual \(\mathrm{Adm2F}\)
input by WRS8.

Apply the complete class-level theorem of LGS13 to these two
\(\mathrm{Closed}\) inputs.  Its coefficient choice can be
written explicitly here: choose exponent cutoff \(1/104\);
for a later \(\epsilon\in(0,1/104]\), take the pointwise maximum
of the zero bounds in WRS7--WRS8 and
\[
 \begin{gathered}
 K_{\mathrm{loc}}.\mathrm{pre}=0,\qquad
 K_{\mathrm{loc}}.t=0,\qquad
 K_{\mathrm{loc}}.\mathrm{mid}=4129,\\
 K_{\mathrm{loc}}.\mathrm{last}
     =\max\{5600002,C_{\mathrm{orb}}(\epsilon)\},\\
 C_{\mathrm{orb}}(\epsilon)
   =40(6r+1)\bigl((h+3)10^{6r+1}+6\bigr),\qquad
 h=2+8\epsilon,\quad r=\lceil8(h+2)/\epsilon\rceil_++1.
 \end{gathered}                                                   \tag{WRS.26}
\]
The bounds are chosen before \(\mathrm{Fm},C_H\).
The given family's \(\mathrm{Large}\) hypothesis at this
bound is exactly the sufficient coefficient region used in LGS13.
Now fix a member and all its constants.  LGS13 and GLG9 choose
one \(C_{\mathrm{count}}\ge0\) and the finitely many eventual
thresholds.  Increase their maximum to include
\(x_g\), (WRS.24), and \(\mathrm{Fm}.x_1+1\).  All later
members \(i\), places \(w\), primes, served representatives,
local degrees and coefficient embeddings use this same threshold.
Set
\[
 C_G=4C_{\mathrm{count}}+1,\qquad
 \delta_G(x)=x^{-\epsilon/2},\qquad
 b_o=(\mathrm{Fm}.L(x,i).\mathrm{sc}.b)^{\delta_G(x)}.
 \tag{WRS.27}
\]
Since \(\epsilon>0\), \(\delta_G(x)\to0\) as \(x\to+\infty\).
The function \(\delta_G\) is chosen once, before any of the
later indices or places.

For precision, the full pointwise conclusion is the literal
\[
 \mathrm{LinkG\_at}(\mathrm{Fm}.G_x(x,i),\mathrm{Fm}.G_K,C_G,b_o).
\]
Write \(D,L,X,K\) for that member's original records and
\(p=p_w\).  Its five conjuncts are the following.
First, for every \(u,w\) with \(p\ne T_u\),
\[
 \begin{gathered}
 \mathrm{LocLe}_{D.\mathrm{plEq}(w)}
    (\lfloor\mathrm{ur}_w(L.\tau(w))\rfloor_+,
          L.\mathrm{NuLat}(u),\mathrm{MramL}(D,u)),\\
 L.\mathrm{lenMN}(u,w)
       \le D.\mathrm{duS}(u)\{L.\tau(w)+L.\mathrm{io}(w)\},\\
 \mathrm{VV}_{p}=\varnothing\Longrightarrow
       L.\tau(w)=L.\mathrm{io}(w)=L.\mathrm{lenMN}(u,w)=0.
 \end{gathered}                                                   \tag{WRS.28}
\]
Second, for every original prime argument and every
\(t_0\in\mathrm{VV}_{p}\),
\[
 S_v(X,K,a_L,b_L,p,t_0)
      \le 2m\,N_v(X,p,t_0)\,L.\mathrm{sc}.a/p.
 \tag{WRS.29}
\]
Third, at every place with the original \(\mathrm{InertK}(D,p)\)
guard and \(p\le L.\mathrm{sc}.X''(D.R.n)\),
\[
 \begin{gathered}
 0\le L.\mathrm{gain}(j,w)
       \le4|TT|L.\mathrm{sc}.a/p
       \qquad(j\in\mathrm{lifts}(D)),\\
 \sum_{j\in\mathrm{lifts}(D)}\ \sum_{t\in L.TT_w}
   \bigl\{L.\mathrm{vonMangoldtClassSum}(j,t,p)
                  -\mathrm{loG}(L,C_G,b_o,p)\bigr\}
       -m\,\mathrm{flG}(L,C_G,p)\\
 \hspace{25mm}\le
       \sum_{j\in\mathrm{lifts}(D)}L.\mathrm{gain}(j,w)-L.\tau(w).
 \end{gathered}                                                   \tag{WRS.30}
\]
Fourth, \(\mathrm{ServedOrbitsContained}(X,K,a_L,b_L)\) holds
for the original orbit representatives.  Fifth, the full original
\(\mathrm{MeanKappaBound}(X,K)\), with both its clauses, holds.

These are exactly the conclusions of LGS13 on the supplied
compatibility and selected tops.  In particular its local length
identification uses the original valuation normalization; it
introduces no extra ramification factor into (WRS.28).  Its
auxiliary zero-chain application leaves the actual \(D,L,X,K\),
comparison, local lattices and all five displayed conclusions
unchanged.  WRS1--WRS7 identified the geometric kernel before
that auxiliary argument; they did not replace it by that zero chain.
Thus the counting, local-length, coverage and mean estimates all
concern the single original member.  Equations (WRS.27)--(WRS.30)
and the last two predicates prove the original \(\mathrm{LinkG}\),
and the preceding coefficient choice proves its full \(\mathrm{Closed}\)
form (WRS.25).
\hfill\(\square\)
% END WRS9

% BEGIN WRS10
\paragraph{WRS10. Literal universal installer and the selected closure boundary.}
\label{claim:WRS10}
The original statement
\[
 \begin{gathered}
 \forall\mathcal A:\mathrm{FamClass},\quad
 \mathrm{TSel}(\mathcal A)\ \Longrightarrow\
 \bigl(\forall\mathrm{Fm}\in\mathcal A,
                         \mathrm{Duty}(\mathrm{Fm})\bigr)
       \ \Longrightarrow\\
 \mathrm{Closed}(D_3(3,\mathcal A),\mathrm{LinkG})
 \end{gathered}                                                   \tag{WRS.31}
\]
is proved with every original quantifier.  The selected built-class
conclusion (WRS.25) is proved separately without either of those
two extra premises.  Both versions retain the same family objects,
the same \(D_3\) conditions and the exact five-field
\(\mathrm{LinkG\_at}\) conclusion.

To prove (WRS.31), fix the arbitrary original class and selector.
At each family member the selector gives one threshold and, after
that threshold, its original \(\mathrm{RecT}\) on every supplied
\(\mathrm{Extra}\).  Its \(T.g\) field is the required
\(\mathrm{GroupCompatibilityAll}\) on that \(\mathrm{Extra}.c\).
GCF17 therefore gives \(\mathrm{GroupCompatibilityFamily}\)
on the same member, with the maximum of the selector threshold
and \(\mathrm{Fm}.x_1+1\).  The elementary \(\mathrm{Closed}\)
argument of WRS7 supplies its first class-level input.
GCF18 supplies the original selected-top input on \(D_3\),
and LGS13 proves the conclusion.  This proof even leaves the
original duty premise unused; retaining it gives exactly
\texttt{GlobalTheory.Top.link\_g}.  The framed version, with
its subsequent repeated \(D_3\)-membership,
\(\mathrm{TuSize}\) and \(\mathrm{DGamSize}\) guards, follows
by ignoring those redundant guards after (WRS.31).

For the actual class selected by the construction, the route is
instead WRS5--WRS9.  Its witnesses are extracted from existing
built-class membership using (WRS.18); no \(\mathrm{TSel}\)
premise is inserted.  If the common construction intersects
\(\mathcal A\) with its target class before adding the original
\(\mathrm{WithChN},\mathrm{WithSm},\mathrm{WithTw}\) and
selected-design predicates, membership still implies membership
in \(\mathcal A\).  WRS7 applies directly to those same objects;
the selected-top projection and the coefficient choice in WRS9
are unchanged.  Equivalently one can restrict the \(\mathrm{Closed}\)
theorem to this fixed subclass, without choosing a new member.

Consequently the original local-gain supplier is complete on the
selected built class.  It does not assert
\(\mathrm{DesignExists}(\mathcal A)\): all statements above are
universal over members, and the original common-design obligation
still constructs a member simultaneously satisfying its other
geometric, row, period and arithmetic requirements.  Nor does
local compatibility manufacture the stronger \(\mathrm{RecT}\)
record needed by the denominator, period, spectral and other
second-layer arguments.  Those distinct original producers remain
separate.  This proves precisely the selected local-gain endpoint,
and makes no assertion of a completed RH or GRH proof.
\hfill\(\square\)
% END WRS10
